\section{Tool Object}

A Tool object describes a cutting tool used by machining operations within a MIFX package.

Tools are defined independently from operations and may be referenced by multiple operations throughout the machining process.

Tool objects are stored within the \texttt{tools} directory and are organized into tool groups.

Example location:

\begin{verbatim}
tools/main/7.json
\end{verbatim}

\subsection{Tool Structure}

A Tool object typically contains:

\begin{itemize}
\item a tool identifier
\item optional descriptive metadata
\item artifact references describing tool properties
\end{itemize}

Example:

\begin{verbatim}
{
  "id": "7",
  "unit": "in",
  "description": "1/2 Bull Mill",

  "artifacts": [
    {
      "role": "tool_profile",
      "kind": "json",
      "path": "profiles/7.json",
      "present": true
    },
    {
      "role": "tool_mount",
      "kind": "json",
      "path": "mounts/7.json",
      "present": true
    },
    {
      "role": "tool_controls",
      "kind": "json",
      "path": "controls/7.json",
      "present": true
    },
    {
      "role": "tool_extensions",
      "kind": "json",
      "path": "extensions/7.json",
      "present": true
    }
  ]
}
\end{verbatim}

\subsection{Tool Identifier}

Each tool must define an identifier using the \texttt{id} field.

Example:

\begin{verbatim}
"id": "7"
\end{verbatim}

Tool numbering and group context are defined by the tool group structure rather than by the filename.

\subsection{Tool Artifacts}

Tool properties are defined through artifacts that reference additional resources.

Common tool artifacts include:

\begin{itemize}
\item tool profile definitions
\item tool mounting configuration
\item control point definitions
\item optional geometry meshes
\item optional extensions
\end{itemize}

These resources are stored as separate files within the tool group directory.

\subsection{Tool Profile}

Tool profiles describe the geometric shape of the cutting tool.

Profiles may be defined using parametric descriptions or geometry files.

Examples include:

\begin{itemize}
\item revolved profiles
\item line and arc segment definitions
\item mesh representations
\end{itemize}

\subsection{Tool Mount}

Tools may optionally define mounting configurations describing how the tool is oriented relative to the machine spindle.

Mount definitions may include transformation matrices representing angle heads or other mounting devices.

\subsection{Tool Controls}

Tool controls describe control system references associated with the tool.

Examples include:

\begin{itemize}
\item length offsets
\item diameter offsets
\item control points used by the machining process
\end{itemize}

\subsection{Tool Extensions}

Additional tool information may be provided using extension artifacts.

Examples may include:

\begin{itemize}
\item holder descriptions
\item vendor information
\item tooling system identifiers
\end{itemize}

Implementations should ignore extensions they do not recognize.
